% Compile with latexmk (uses ~/.latexmkrc: LuaLaTeX, output in build/)
% Read-aloud teaching script for Week 8 (companion to beamer.tex, follows it
% frame by frame). Sessions 1 and 2 are held on the SAME DAY, back to back
% (2 x 90 min with a short break in between).
% Exercises are numbered by week (8.1, 8.2, ...); finance examples share the
% numbering. Each one lives in exercises/NN-slug.tex.

\documentclass[11pt,a4paper,parskip=half]{scrartcl}
\usepackage{amsmath}
\usepackage[stixtwo]{fontsetup}
\usepackage[english]{babel}
\usepackage{booktabs}
\usepackage{enumitem}
\usepackage{tikz}
\usepackage[most]{tcolorbox}
\usepackage[a4paper,margin=1in,marginparwidth=1.9cm,marginparsep=2mm]{geometry}
\usepackage[hidelinks]{hyperref}

\definecolor{navy}{RGB}{24,48,75}
\definecolor{teal}{RGB}{0,120,125}
\definecolor{gold}{RGB}{190,130,20}
\definecolor{plum}{RGB}{110,60,120}

\addtokomafont{disposition}{\color{navy}}
\setkomafont{section}{\Large\bfseries\color{navy}}

% Contact
\newcommand{\myemail}{tiborsandor2001@gmail.com}

\newcommand{\Om}{\Omega}

% Exercise statement (what students see on the slide)
\newtcolorbox{task}[1]{enhanced,breakable,colback=navy!4,colframe=navy,
  fonttitle=\bfseries,title={#1}}
% Finance example
\newtcolorbox{fin}[1]{enhanced,breakable,colback=teal!5,colframe=teal,
  fonttitle=\bfseries,title={#1}}
% Teaching cue, not read aloud
\newtcolorbox{cue}{enhanced,breakable,colback=gold!8,colframe=gold,
  boxrule=.4pt,left=4pt,right=4pt,top=2pt,bottom=2pt,fontupper=\small\itshape}
% Background for the lecturer (finance theory), NOT read aloud
\newtcolorbox{bg}[1]{enhanced,breakable,colback=plum!4,colframe=plum,
  boxrule=.6pt,fonttitle=\bfseries\small,
  title={Background for you (not read aloud): #1},fontupper=\small,
  before upper={\setlength{\parskip}{4pt}}}
% Print version (script-print.tex defines \scriptprint): the background
% boxes are dropped entirely.
\ifdefined\scriptprint
  \RenewDocumentEnvironment{bg}{m +b}{}{}
\fi
\newcommand{\terms}{\par\textbf{\color{plum}Key terms.}\ }
\newcommand{\sq}[1]{\par\textbf{\color{plum}Student may ask:}\ \emph{#1}\par}
\newcommand{\ans}{\textbf{Answer:}\ }

\newcommand{\say}{\textbf{\color{teal}Say:}\ }
\newcommand{\board}{\textbf{\color{gold}Board:}\ }
\newcommand{\ask}{\textbf{\color{navy}Ask the class:}\ }
\newcommand{\slide}[1]{\marginpar{\raggedright\footnotesize\color{gray}slide #1}}

% Small Venn diagram of two events E, F in Omega (same as in the deck)
\newcommand{\venn}[1]{%
  \begin{tikzpicture}[scale=.5,baseline=(O.base)]
    \node (O) at (.5,0) {\phantom{x}};
    #1
    \draw[thick] (-1.5,-1.15) rectangle (2.5,1.15);
    \draw[thick,navy] (0,0) circle (.85);
    \draw[thick,navy] (1,0) circle (.85);
    \node at (-.35,0) {\scriptsize $E$};
    \node at (1.35,0) {\scriptsize $F$};
    \node at (2.15,.8) {\scriptsize $\Omega$};
  \end{tikzpicture}}
\newcommand{\shadeall}[1]{\fill[#1] (-1.5,-1.15) rectangle (2.5,1.15);}

% The exercises (one file each in exercises/) and \scriptexercise
% Exercises of Week 8, shared by beamer.tex, the study guide, the script and
% exercises.tex. Every exercise lives in a file of its own,
% exercises/NN-slug.tex, numbered 8.NN (see exercises.lua for the docstring).
% The body of an exercise file:
%
%   \begin{problem} ... \end{problem}      the statement, used everywhere;
%                                          inside it, \begin{context} ...
%                                          \end{context} is background that
%                                          the slides leave out
%   \begin{slideproblem} ... \end{slideproblem}
%                                          (optional) a shorter statement
%                                          for the slides
%   \begin{slides} ... \end{slides}        one frame of the deck: \exstatement
%                                          places the statement; repeat the
%                                          environment for more frames
%   \begin{solution} ... \end{solution}    the detailed solution (study
%                                          guide, exercise PDFs, website)
%
% Commands for the documents (the argument is the slug: 01-license-plates.tex
% has the slug license-plates):
%
%   \exref{slug}                 the number, e.g. 8.1, for "Exercise~\exref{...}"
%   \extitle{slug}               the title from the docstring
%   \exerciseframe[opts]{slug}   beamer: the frames of the exercise
%   \studyexercise{slug}         study guide: heading, statement, solution
%   \scriptexercise{slug}        script: the statement box, as on the slide
%   \allexercises                exercises.tex: every exercise, in order; on
%                                the website statement and solution with
%                                markers where the pipeline splits the HTML
%                                (\collectionexercise), in the PDF plain:
%                                each statement, followed by its solution if
%                                \ifsolutions (exercises-solutions.tex)

\def\exweek{8}

\directlua{exercises = dofile("exercises/exercises.lua")
  exercises.scan("exercises", "\exweek")}

\ExplSyntaxOn

\msg_new:nnn { exercises } { unknown } { Unknown~exercise~`#1'. }

\cs_new_protected:Npn \__ex_missing:n #1
  { \msg_error:nnn { exercises } { unknown } {#1} ?? }
\cs_new:Npn \exref #1
  {
    \cs_if_exist:cTF { exnum@#1 }
      { \use:c { exnum@#1 } } { \__ex_missing:n {#1} }
  }
\cs_new:Npn \extitle #1
  {
    \cs_if_exist:cTF { extitle@#1 }
      { \use:c { extitle@#1 } } { \__ex_missing:n {#1} }
  }

% ---------------------------------------------------------------------------
% Loading an exercise: its parts are stored, nothing is typeset. The slides
% typeset the stored parts; the other documents read the part they need from
% the file again (\__ex_input_part:nn), as stored tables do not survive lwarp.

\tl_new:N \g__ex_current_tl
\tl_new:N \g__ex_problem_tl
\tl_new:N \g__ex_slideproblem_tl
\tl_new:N \g__ex_solution_tl
\seq_new:N \g__ex_slides_seq
\box_new:N \l__ex_scratch_box

\cs_new_protected:Npn \__ex_load:n #1
  {
    \cs_if_exist:cF { exfile@#1 } { \__ex_missing:n {#1} }
    \tl_gset:Nn \g__ex_current_tl {#1}
    \tl_gclear:N \g__ex_problem_tl
    \tl_gclear:N \g__ex_slideproblem_tl
    \tl_gclear:N \g__ex_solution_tl
    \seq_gclear:N \g__ex_slides_seq
    \vbox_set:Nn \l__ex_scratch_box
      {
        \DeclareDocumentEnvironment { problem } { +b }
          { \tl_gset:Nn \g__ex_problem_tl {##1} } { }
        \DeclareDocumentEnvironment { slideproblem } { +b }
          { \tl_gset:Nn \g__ex_slideproblem_tl {##1} } { }
        \DeclareDocumentEnvironment { slides } { +b }
          { \seq_gput_right:Nn \g__ex_slides_seq {##1} } { }
        \DeclareDocumentEnvironment { solution } { +b }
          { \tl_gset:Nn \g__ex_solution_tl {##1} } { }
        \exp_args:Nv \input { exfile@#1 }
      }
  }

% Typesets one part (problem, slideproblem, solution) of an exercise straight
% from its file; the other parts are skipped.
\cs_new_protected:Npn \__ex_input_part:nn #1#2
  {
    \group_begin:
      \clist_map_inline:nn { problem, slideproblem, slides, solution }
        {
          \str_if_eq:nnTF {##1} {#2}
            { \DeclareDocumentEnvironment {##1} { } { } { } }
            { \DeclareDocumentEnvironment {##1} { +b } { } { } }
        }
      \exp_args:Nv \input { exfile@#1 }
    \group_end:
  }

% Exercise 8.1, or Exercise 8.2 · Finance
\cs_new:Npn \__ex_heading:
  {
    Exercise~\exref { \g__ex_current_tl }
    \cs_if_exist:cT { exfinance@ \g__ex_current_tl } { ~·~Finance }
  }
\prg_new_conditional:Npnn \__ex_if_finance: { T, F, TF }
  {
    \cs_if_exist:cTF { exfinance@ \g__ex_current_tl }
      { \prg_return_true: } { \prg_return_false: }
  }

% ---------------------------------------------------------------------------
% Beamer

% The finance context is printed in the study guide, the exercise PDFs and on
% the website; the
% slides and the script (which shows what the slide shows) leave it out.
\bool_new:N \l__ex_context_bool
\bool_set_true:N \l__ex_context_bool
\DeclareDocumentEnvironment { context } { }
  {
    \bool_if:NTF \l__ex_context_bool
      { \par \textit { Finance~context. } ~ }
      { \__ex_gobble_context:w }
  }
  { \bool_if:NT \l__ex_context_bool { \par } }
\cs_new_protected:Npn \__ex_gobble_context:w #1 \end #2
  { \str_if_eq:nnTF {#2} { context } { \end { context } } { \__ex_gobble_context:w } }

\IfClassLoadedTF { beamer }
  {
    \bool_set_false:N \l__ex_context_bool

    \cs_new:Npn \__ex_block: { \__ex_if_finance:TF { exampleblock } { block } }
    \NewDocumentCommand \exstatement { }
      {
        \begin { \__ex_block: } { \__ex_heading: }
          \tl_if_empty:NTF \g__ex_slideproblem_tl
            { \g__ex_problem_tl } { \g__ex_slideproblem_tl }
        \end { \__ex_block: }
      }

    \tl_new:N \l__ex_frames_tl
    \NewDocumentCommand \exerciseframe { O{t} m }
      {
        \__ex_load:n {#2}
        \tl_clear:N \l__ex_frames_tl
        \seq_map_indexed_inline:Nn \g__ex_slides_seq
          {
            \tl_put_right:Nn \l__ex_frames_tl { \begin{frame}[#1] }
            \tl_put_right:Ne \l__ex_frames_tl
              {
                {
                  \exp_not:N \extitle {#2}
                  \int_compare:nNnT {##1} > { 1 } { \c_space_tl (cont.) }
                }
              }
            \tl_put_right:Nn \l__ex_frames_tl { ##2 \end{frame} }
          }
        \tl_use:N \l__ex_frames_tl
      }
  }
  {
    \cs_new:Npn \__ex_box: { \__ex_if_finance:TF { fin } { task } }
    \cs_new_protected:Npn \__ex_statement_box:n #1
      {
        \begin { \__ex_box: } { \__ex_heading: }
          \__ex_input_part:nn { \g__ex_current_tl } {#1}
        \end { \__ex_box: }
      }

    \NewDocumentCommand \studyexercise { m }
      {
        \__ex_load:n {#1}
        \subsection* { Exercise~\exref{#1}~---~\extitle{#1} }
        \addcontentsline { toc } { subsection }
          { Exercise~\exref{#1}~---~\extitle{#1} }
        \__ex_statement_box:n { problem }
        \soltitle \__ex_input_part:nn {#1} { solution }
      }

    \NewDocumentCommand \scriptexercise { m }
      {
        \__ex_load:n {#1}
        \group_begin:
          \bool_set_false:N \l__ex_context_bool
          \tl_if_empty:NTF \g__ex_slideproblem_tl
            { \__ex_statement_box:n { problem } }
            { \__ex_statement_box:n { slideproblem } }
        \group_end:
      }

    % The website (lwarp): statement and solution of each exercise, between
    % markers where the pipeline splits the HTML.
    \NewDocumentCommand \collectionexercise { m }
      {
        \__ex_load:n {#1}
        \par [[exercise:#1:problem]] \par
        \__ex_input_part:nn {#1} { problem }
        \par [[exercise:#1:solution]] \par
        \__ex_input_part:nn {#1} { solution }
        \par [[exercise:#1:end]] \par
      }
    % The PDF: plain, each statement, followed by its solution in the
    % version with solutions (\ifsolutions). The statement is set off like a
    % quote in Markdown, by a bar on the left (quote=<color>, also used by
    % exercises.tex) in the color of its box in the study guide: navy, or teal
    % for a finance exercise.
    \tcbset
      {
        quote/.style =
          {
            enhanced, breakable, frame~hidden, interior~hidden,
            sharp~corners, boxrule = 0pt, borderline~west = {2pt}{0pt}{#1},
            left = 8pt, right = 0pt, top = 1pt, bottom = 1pt,
            before~skip = 6pt, after~skip = 6pt,
            % the document's paragraph spacing (parskip=half)
            before~upper = { \setlength { \parskip } { .5\baselineskip } }
          }
      }
    \NewDocumentCommand \plainexercise { m }
      {
        \__ex_load:n {#1}
        \subsection* { Exercise~\exref{#1}~---~\extitle{#1} }
        \addcontentsline { toc } { subsection }
          { Exercise~\exref{#1}~---~\extitle{#1} }
        \__ex_if_finance:TF
          { \begin{tcolorbox}[quote=teal] }
          { \begin{tcolorbox}[quote=navy] }
          \__ex_input_part:nn {#1} { problem }
        \end{tcolorbox}
        \legacy_if:nT { solutions }
          {
            \soltitle
            \__ex_input_part:nn {#1} { solution }
          }
      }
    \NewDocumentCommand \allexercises { }
      {
        \IfPackageLoadedTF { lwarp }
          { \directlua { exercises.each("collectionexercise") } }
          { \directlua { exercises.each("plainexercise") } }
      }
  }

\ExplSyntaxOff


% PDF properties and the date on the title page
% The properties of every PDF of Week 8: author (with email address), title,
% subject, keywords and the date the document last changed, in the PDF
% (hyperref) and on its title page (\date), and the header and footer of the
% printed documents (below). Each document loads this file after hyperref and
% then calls \documentinfo{<what it is>}, e.g. \documentinfo{Slides}.
%
% The date: the website build (scripts/build-content.mjs) puts
% \def\lastmodified{YYYY-MM-DD} in front of the document, the date of the last
% commit that changed its sources. Compiled locally, it is today.

\newcommand{\courseinfo}{Mathematics for Finance}
\newcommand{\weeklabel}{Week 8}
\newcommand{\weekinfo}{\weeklabel: Combinatorial analysis, sample space, events}
\newcommand{\authorname}{Tibor SÁNDOR}
\providecommand{\myemail}{tiborsandor2001@gmail.com}

\makeatletter
\ifdefined\lastmodified\else
	\edef\lastmodified{\the\year-\two@digits{\the\month}-\two@digits{\the\day}}
\fi
\def\di@split#1-#2-#3\@nil{%
	\def\di@year{#1}\def\di@month{#2}\def\di@day{#3}}
\expandafter\di@split\lastmodified\@nil
% November 11, 2026
\edef\lastmodifiedtext{%
	\ifcase\numexpr\di@month\relax\or January\or February\or March\or April\or
		May\or June\or July\or August\or September\or October\or November\or
		December\fi\space\number\di@day, \di@year}
% D:20261111000000, the date format of PDF metadata
\edef\di@pdfdate{D:\di@year\di@month\di@day000000}

\newcommand{\documentinfo}[1]{%
	\hypersetup{
		pdftitle={\courseinfo, \weekinfo\space -- #1},
		pdfauthor={\authorname\space<\myemail>},
		pdfsubject={\courseinfo, \weekinfo},
		pdfkeywords={probability, combinatorics, mathematics for finance,
				Corvinus University of Budapest},
		pdfcreationdate={\di@pdfdate},
		pdfmoddate={\di@pdfdate},
		pdflang={en-US},
	}%
	\date{Last updated: \lastmodifiedtext}%
	\gdef\documentkind{#1}%
}

% Header and footer of the printed documents (KOMA-Script, scrlayer-scrpage;
% the slides have beamer's footline, the website neither): the course, the
% week and the document on the left of the header (the week's topic is on the
% title page), the university on the right; in the footer the author and the
% page as <current>/<total>. The title page gets only the footer.
\newcommand{\documentkind}{}
\IfClassLoadedTF{beamer}{}{\IfPackageLoadedTF{lwarp}{}{%
	\usepackage{lastpage}%
	\usepackage[headsepline=1.5pt,footsepline=1.5pt,plainfootsepline]
		{scrlayer-scrpage}%
	\clearpairofpagestyles
	\setkomafont{pageheadfoot}{\small\sffamily\color{navy}}%
	\setkomafont{pagenumber}{}%
	\setkomafont{headsepline}{\color{navy!35}}%
	\setkomafont{footsepline}{\color{navy!35}}%
	\ihead{\courseinfo\ \textperiodcentered\ \weeklabel\ \textperiodcentered\ 
		\documentkind}%
	\ohead{Corvinus University of Budapest}%
	\ifoot*{\authorname\ \textperiodcentered\ \myemail}%
	\ofoot*{\pagemark/\pageref*{LastPage}}%
	\pagestyle{scrheadings}%
	% Extra stretch for breaking lines with long inline formulas.
	\setlength{\emergencystretch}{1.5em}%
}}
\makeatother

\documentinfo{Teaching script\ifdefined\scriptprint\space(print version)\fi}

\title{Mathematics for Finance}
\subtitle{Week 8 --- Combinatorial analysis, sample space, events\\[2pt]
  \normalsize Teaching script for two 90-minute practical classes held on the
  same day, back to back\ifdefined\scriptprint\\[2pt]
  (print version, without the background boxes)\fi}
\author{Tibor SÁNDOR\\[2pt]
  \normalsize\href{mailto:\myemail}{\texttt{\myemail}}}

\begin{document}
\maketitle

\begin{cue}
	How to use this script. Plain text is meant to be read out loud or
	paraphrased. \textbf{Say} marks key sentences; \textbf{Ask} marks questions
	to put to the class (wait for answers before revealing the solution);
	\textbf{Board} marks what to write. Gold boxes like this one are short notes
	for you only.\ifdefined\scriptprint\else\ \textbf{Purple ``Background for
			you'' boxes} explain the finance behind the examples (what the terms mean,
		why the example matters, what students may ask); they are \emph{not} meant
		to be read out, but they let you answer follow-up questions with
		confidence.\fi\ The script follows the
	deck frame by frame; the slide numbers in the margin are the frame numbers
	in the footline of the deck. Exercise numbers match the slides.
	Timings are approximate; items marked \emph{(optional)} can be dropped if
	you run late.
\end{cue}

\subsection*{Timing overview}

Both sessions are on the same day: Session 1 (90 min), a short break,
Session 2 (90 min).

\begin{center}
	\small
	\begin{tabular}{@{}llrr@{}}
		\toprule
		Slides & Session 1                                        & min & cum.        \\\midrule
		1--3   & Introduction (who I am, contact), plan for today & 6   & 6           \\
		4      & Toolbox                                          & 6   & 12          \\
		5      & Exercise~\exref{license-plates}                  & 3   & 15          \\
		6      & Exercise~\exref{pin-codes} (PIN codes)           & 3   & 18          \\
		7      & Exercise~\exref{couples-at-a-concert}            & 9   & 27          \\
		8      & Exercise~\exref{bus-queue}                       & 8   & 35          \\
		9      & Exercise~\exref{letter-arrangements}             & 7   & 42          \\
		10     & Exercise~\exref{eight-digit-numbers}             & 5   & 47          \\
		11     & Exercise~\exref{up-and-down-days} (up and down days) & 3   & 50          \\
		12     & Exercise~\exref{awards}                          & 3   & 53          \\
		13     & Exercise~\exref{ranking-funds} (ranking mutual funds) & 4   & 57          \\
		14     & Exercise~\exref{committees}                      & 10  & 67          \\
		15     & Exercise~\exref{portfolio} (portfolio), lottery  & 7   & 74          \\
		16     & Exercise~\exref{dance-class} (optional)          & 3   & 77          \\
		17     & Exercise~\exref{lattice-paths} (lattice paths)   & 9   & 86          \\
		18     & Summary; short break                             & 4   & 90          \\\midrule
		       & \textbf{Total}                                   &     & \textbf{90} \\\bottomrule
	\end{tabular}
\end{center}

\begin{center}
	\small
	\begin{tabular}{@{}llrr@{}}
		\toprule
		Slides & Session 2 (after the break)       & min & cum.        \\\midrule
		19--20 & After the break: two-minute recap & 3   & 3           \\
		21     & Toolbox                           & 5   & 8           \\
		22     & Exercise~\exref{bridge}          & 5   & 13          \\
		23     & Exercise~\exref{dividing-into-groups} & 4   & 17          \\
		24     & Exercise~\exref{rating-bonds} (rating bonds) & 2   & 19          \\
		25     & Exercise~\exref{expansions}, identities & 6   & 25          \\
		26     & Exercise~\exref{stars-and-bars}  & 4   & 29          \\
		27     & Exercise~\exref{identical-dice}  & 2   & 31          \\
		28     & Exercise~\exref{allocating-capital} & 5   & 36          \\
		29     & Exercise~\exref{pascals-rule}    & 3   & 39          \\
		30     & Exercise~\exref{binomial-tree} (binomial tree) & 5   & 44          \\
		31     & Pascal's triangle                 & 4   & 48          \\
		32     & Sample space and events           & 5   & 53          \\
		33     & De Morgan's laws for two events   & 5   & 58          \\
		34     & De Morgan's laws: general form    & 2   & 60          \\
		35     & Exercise~\exref{sample-space}    & 3   & 63          \\
		36     & Exercise~\exref{stock-path-events} & 4   & 67          \\
		37     & Exercise~\exref{two-dice}        & 5   & 72          \\
		38     & Exercise~\exref{credit-basket}   & 5   & 77          \\
		39--40 & Exercise~\exref{newspapers}      & 6   & 83          \\
		41     & Exercise~\exref{bank-products}   & 4   & 87          \\
		42     & Summary                           & 3   & 90          \\\midrule
		       & \textbf{Total}                    &     & \textbf{90} \\\bottomrule
	\end{tabular}
\end{center}

\tableofcontents

%======================================================================
\section{Session 1: Permutations and combinations}
%======================================================================

\subsection{Introduction and plan for today (6 min)}\slide{1--3}

\begin{cue}
	Slide 1 (title) is on the screen while you introduce yourself. Speak
	slowly; this is the first time they see you.
\end{cue}

Good morning, everyone, and welcome. Before we start, let me introduce
myself. My name is Tibor Sándor --- on the slides you see SÁNDOR in capitals,
that is my family name. I am new to this course: from today I will be
teaching you the probability part of Mathematics for Finance.

My background is in engineering, not in finance. I actually think that is
an advantage for this part of the course. Engineers use probability every
day: how likely is it that a system fails, how large are measurement errors,
how much safety margin do we need. You will see that a lot of finance is
the same mathematics with different names --- a bond that defaults behaves
very much like a component that fails, and diversifying a portfolio is very
much like building redundancy into a system. So I will show you the
mathematics carefully, and I will always tell you the financial story behind
the examples. And if some of you know more about a financial product than I
do, please speak up --- I am happy to learn from you as well.

How to reach me: my email address is
\href{mailto:\myemail}{\texttt{\myemail}}. It is on the title slide and on
the next slide. Please put ``Math for Finance'' in the subject line, so that
I find your message quickly. You can also just catch me after class.

\begin{cue}
	Slide 2 (Plan for today).
\end{cue}

\say Today is a long day: we have two sessions back to back, 90 minutes
each, with a short break in between.

In Session 1 we count: the counting principle, permutations, permutations
with repetition, combinations and lattice paths. After the break, in
Session 2, we divide objects into groups, we meet the binomial and
multinomial theorems and the ``stars and bars'' method, we see how counting
works in the binomial tree of a stock price and arrange the counts in
Pascal's triangle, and then we introduce the language of probability: the
sample space, events, De Morgan's laws and Venn diagrams.

Why start the probability part with counting? Because when all outcomes are
equally likely, a probability is simply (favorable outcomes)/(all outcomes)
--- and that is exactly where we start next week.

\begin{cue}
	Slide 3 (Session 1 title) --- just click through.
\end{cue}

\begin{bg}{the big picture of ``finance and counting''}
	You do not need finance knowledge to teach this week, but here is the
	vocabulary that appears in the examples, in one place.

	\emph{Stock (share)}: a piece of ownership in a company; its price moves
	randomly. \emph{Bond}: a loan to a company or a government; the borrower
	(the \emph{issuer}) pays interest and repays the face value at maturity.
	If the issuer cannot pay, the bond \emph{defaults} --- in engineering terms,
	the component fails. \emph{Portfolio}: a collection of investments held
	together. \emph{Derivative}: a contract whose payoff depends on something
	else (a stock price, whether a bond defaults, \ldots). The two derivatives
	in this week's examples are \emph{options} (the right, but not the
	obligation, to buy or sell a stock at a fixed price later) and \emph{credit
		default swaps} (insurance against a bond defaulting).

	Why counting matters: many pricing and risk models are built on a finite
	set of scenarios. If all scenarios are equally likely, probability is
	counting; if they are not (as in the binomial model with $q\neq 1/2$), we
	still need to know \emph{how many} scenarios lead to the same outcome,
	because they all get the same weight. That is exactly what binomial
	coefficients do in the binomial tree.

	\terms stock (share); bond; issuer; default; portfolio; derivative;
	option; credit default swap; scenario.

	\sq{Is finance mathematics really the same as engineering mathematics?}
	\ans The tools are the same (probability, statistics, optimization,
	differential equations); the difference is that prices are set by people
	who react to each other, so models are less stable than physical laws. A
	good engineer's habit --- ``check the model against reality and know its
	limits'' --- is exactly what finance needs.
\end{bg}

\subsection{Toolbox (6 min)}\slide{4}

\say The \emph{basic principle of counting}: if something can be done in
$m$ ways, and then, for each of those, a second thing in $n$ ways, then the
two together can be done in $m\cdot n$ ways. With more steps we keep
multiplying: $n_1 n_2\cdots n_r$.

\board $m$ ways, then $n$ ways $\Rightarrow$ $m\cdot n$ ways.

Think of a tree: the first step gives $m$ branches, each branch splits
into $n$ branches, so there are $m\cdot n$ leaves. From this single principle
we get everything else.

\emph{Permutations --- order matters.} In how many ways can we arrange $n$
distinct objects in a row? The first place can be filled in $n$ ways, the
second in $n-1$ ways (one object is used up), and so on, so $n!$. If we only
fill $k$ places, we stop after $k$ factors:
\[ \frac{n!}{(n-k)!} = n(n-1)\cdots(n-k+1). \]

\emph{Combinations --- order does not matter.} If we only care \emph{which}
$k$ objects are chosen, look at the ordered selections above: each set of
$k$ objects comes in $k!$ orders, so it was counted $k!$ times. We divide:
\[ \binom nk = \frac{n!}{k!\,(n-k)!}. \]

\begin{cue}
	Permutations with repetition (identical objects, divide by
	$n_1!\,n_2!\cdots n_r!$) are not on the toolbox slide; they come with
	Exercise~\exref{letter-arrangements} (slide 9), where the formula is derived on the MISSISSIPPI
	example.
\end{cue}

\say Before every exercise, ask yourself two questions. First: does the order
matter? Second: can the same element be used more than once? The answers tell
you which formula to use.

\board
\begin{tabular}{lcc}
	\toprule
	                   & order matters & order does not matter          \\\midrule
	no repetition      & $n!/(n-k)!$   & $\binom nk$                    \\
	repetition allowed & $n^k$         & $\binom{n+k-1}{k}$ (Session 2) \\\bottomrule
\end{tabular}

\begin{cue}
	Leave this table on the board for the whole day: you will fill in the
	bottom-right cell after the break, in Session 2.
\end{cue}

\subsection{Exercise~\exref{license-plates} --- License plates (3 min)}\slide{5}

\scriptexercise{license-plates}

\ask Does the order matter? \emph{(Yes --- AB123 and BA123 are different
	plates.)} Can we repeat? \emph{(In (a) yes, in (b) no.)}

\say We fill the 7 places one by one and multiply --- that is the basic
principle.

(a) Repetition is allowed: each letter place has 26 choices and each digit
place has 10 choices:
\[ 26^2\cdot 10^5 = 676\cdot 100\,000 = 67\,600\,000. \]

(b) Now every choice uses up a symbol. Letters: $26\cdot25 = 650$. Digits:
$10\cdot9\cdot8\cdot7\cdot6 = 30\,240$. Together
\[ 650\cdot 30\,240 = 19\,656\,000. \]

\subsection{Exercise~\exref{pin-codes} --- PIN codes, finance (3 min)}\slide{6}

\scriptexercise{pin-codes}

\ask Same two questions: does the order matter, can a digit repeat?
\emph{(Order matters --- 1234 and 4321 are different PINs; repetition is
	allowed in (a), not in (b).)}

(a) Four places, 10 digits each: $10^4 = 10\,000$ PINs.

(b) No repeated digit:
\[ 10\cdot9\cdot8\cdot7 = \frac{10!}{6!} = 5040. \]

(c) Three different guesses cover 3 of the $10\,000$ PINs: the fraction
$3/10\,000$.

\say Next week this fraction becomes a probability: the chance that a thief
who tries 3 random PINs before the card is blocked succeeds. For now we only
need the denominator.

\begin{bg}{PIN security}
	\emph{Intuition.} Security of a PIN is a counting argument: an attacker who
	knows nothing must guess, and the number of possible PINs is the size of
	the haystack. Banks keep the haystack small enough to remember (4 digits)
	and protect it with an \emph{attempt limit}: after typically 3 wrong tries
	the card is blocked or retained. Hence the success chance of blind guessing
	is $3/10\,000 = 0.03\%$.

	\emph{The catch: people do not choose uniformly.} Studies of leaked
	4-digit codes found that a handful of codes (1234, 1111, 0000, birth years
	such as 19xx) account for a large share of all codes; 1234 alone is
	reported at roughly one in ten. So the real guessing probability is much
	higher than $3/10\,000$ --- the ``equally likely'' assumption fails. This is
	a nice preview of next week: classical probability is only as good as the
	equally-likely assumption.

	\terms brute-force attack; attempt limit (lockout); entropy (in bits,
	$\log_2 10\,000\approx 13.3$).

	\sq{Would banning repeated digits make PINs safer?}
	\ans No: it removes weak PINs like 1111, but it shrinks the space from
	$10\,000$ to $5040$, so blind guessing becomes almost twice as easy. Every
	rule that restricts choices reduces the count.

	\sq{Why not 6-digit PINs?}
	\ans Some banks use them ($10^6$ possibilities); it is a trade-off between
	security and how easily customers remember them.
\end{bg}

\subsection{Exercise~\exref{couples-at-a-concert} --- Five couples at a concert (9 min)}\slide{7}

\scriptexercise{couples-at-a-concert}

(a) Ten distinct people in ten seats: $10! = 3\,628\,800$.

(b) \say When some people must sit together, glue them into one block.

We have 5 blocks (the couples). The blocks can be arranged in $5! = 120$
ways. Inside each block the husband can sit on the left or on the right: 2
ways, and this happens independently for all 5 couples, giving $2^5 = 32$. So
\[ 5!\cdot 2^5 = 120\cdot 32 = 3840. \]

(c) Now there are two blocks: the men and the women. The two blocks can be
ordered in $2$ ways (men on the left or women on the left). Inside the men's
block, $5!$ orders; inside the women's block, $5!$ orders:
\[ 2\cdot5!\cdot5! = 2\cdot120\cdot120 = 28\,800. \]

(d) If they alternate, the pattern is either MWMWMWMWMW or WMWMWMWMWM: 2
patterns. Then the men fill their 5 seats in $5!$ ways and the women theirs in
$5!$ ways: again $2\cdot5!\cdot5! = 28\,800$.

\ask Is it a coincidence that (c) and (d) give the same answer? \emph{(No: in
	both cases we choose one of two ``seat patterns'' for the genders, and then
	permute men and women within their seats.)}

\board \textbf{Trick:} ``together'' $\Rightarrow$ glue into a block, permute
the blocks, then permute inside each block.

\subsection{Exercise~\exref{bus-queue} --- Queueing for the bus (8 min)}\slide{8}

\scriptexercise{bus-queue}

(a) $8! = 40\,320$.

(b) Glue the three friends into a block. Now we arrange 6 units (the block and
5 other people): $6! = 720$. Inside the block: $3! = 6$. So
$6\cdot 720 = 4320$.

(c) \say ``Not next to each other'' is awkward to count directly. Count the
opposite --- the bad lines --- and subtract.

The bad lines are those in which the two people \emph{are} adjacent: glue
them into a block, $7$ units, $7! = 5040$ orders, and $2!$ orders inside the
block: $2\cdot 5040 = 10\,080$. Hence
\[ 8! - 2!\cdot 7! = 40\,320-10\,080 = 30\,240. \]

\board Complement principle: $\#\text{good} = \#\text{all} - \#\text{bad}$.

\begin{cue}
	Stress this: the complement trick is the single most useful technique in
	this chapter and will return constantly next week as $P(E^c)=1-P(E)$.
\end{cue}

\subsection{Exercise~\exref{letter-arrangements} --- Letter arrangements (7 min)}\slide{9}

\scriptexercise{letter-arrangements}

\say This is the one formula that is not on the toolbox slide:
\emph{permutations with repetition}. If all letters were different, we would
have $n!$ arrangements. Swapping identical letters does not change the word, so every word is
counted $n_1!\,n_2!\cdots$ times: we divide by the factorials of the
multiplicities.

(a) Five different letters: $5! = 120$.

(b) PROPOSE has 7 letters; P appears twice and O appears twice. If we
labeled them $P_1, P_2, O_1, O_2$ we would get $7!$ arrangements. But
swapping $P_1$ and $P_2$ gives the same word, and so does swapping $O_1$ and
$O_2$. Therefore each word was counted $2!\cdot2!$ times:
\[ \frac{7!}{2!\,2!} = \frac{5040}{4} = 1260. \]

(c) MISSISSIPPI: 11 letters, M once, I four times, S four times, P twice.
Each word is counted $4!\,4!\,2! = 24\cdot24\cdot2 = 1152$ times among the
$11!$ orders:
\[ \frac{11!}{4!\,4!\,2!} = \frac{39\,916\,800}{1152} = 34\,650. \]

(d) ARRANGE: A twice, R twice, N, G, E once: $\dfrac{7!}{2!\,2!} = 1260$.

\board Permutations with repetition: $n$ objects, $n_1$ of one kind, $n_2$
of another, \ldots:
\[ \frac{n!}{n_1!\,n_2!\cdots n_r!}. \]

\subsection{Exercise~\exref{eight-digit-numbers} --- Eight-digit numbers (5 min)}\slide{10}

\scriptexercise{eight-digit-numbers}

\ask What is the catch? \emph{(A number cannot start with 0.)}

All arrangements of the cards: 8 cards, digit 1 twice, digit 3 three times:
\[ \frac{8!}{2!\,3!} = \frac{40\,320}{12} = 3360. \]

The bad ones start with 0 --- they are not 8-digit numbers. Fix the 0 in
front; the remaining 7 cards $1,1,2,3,3,3,4$ can be arranged in
$\dfrac{7!}{2!\,3!} = \dfrac{5040}{12} = 420$ ways. So the answer is
\[ \frac{8!}{2!\,3!}-\frac{7!}{2!\,3!} = 3360 - 420 = 2940. \]

\subsection{Exercise~\exref{up-and-down-days} --- Up and down days, finance (3 min)}\slide{11}

\scriptexercise{up-and-down-days}

\ask Which exercise is this, in disguise? \emph{(Letter arrangements: a
	word with 7 U's and 5 D's.)}

This is exactly a word with 7 U's and 5 D's:
\[ \frac{12!}{7!\,5!} = \binom{12}{5} = 792. \]

\say Notice that permutations with two kinds of objects are the same as
combinations: we just choose \emph{which} 5 of the 12 days were down-days.

\begin{bg}{up/down sequences and random walks}
	\emph{Intuition.} Recording only ``up'' or ``down'' for each trading day
	turns a price chart into a binary word. The simplest model of stock prices,
	the \emph{random walk}, says each day is up or down independently, like a
	coin. Under that model all $2^{12} = 4096$ sequences of 12 days are equally
	likely, and the 792 sequences with exactly 7 ups are all equally likely too
	--- so a pattern such as ``7 ups in a row, then 5 downs'' is no more
	special than any other ordering. Chart-readers (``technical analysts'')
	look for patterns; statisticians test whether the patterns are more
	frequent than counting predicts (e.g.\ the \emph{runs test}).

	\terms up-day / down-day; random walk; streak or run; technical analysis.

	\sq{If there were 7 up-days and only 5 down-days, did the stock rise?}
	\ans Not necessarily: the count ignores the \emph{size} of the moves. Five
	big falls can outweigh seven small rises. In the binomial tree (Exercise~\exref{binomial-tree},
	after the break) we fix the sizes ($u$ and $d$), and then the
	count of ups does determine the final price.
\end{bg}

\subsection{Exercise~\exref{awards} --- Awards, with or without repetition (3 min)}\slide{12}

\scriptexercise{awards}

\say Hand out the awards one by one.

(a) Every award can go to any of the 30 students: repetition is allowed, and
order matters because the awards are different. $30^5 = 24\,300\,000$.

(b) After each award one student fewer is eligible:
\[ 30\cdot29\cdot28\cdot27\cdot26 = \frac{30!}{25!} = 17\,100\,720. \]

\subsection{Exercise~\exref{ranking-funds} --- Ranking mutual funds, finance (4 min)}\slide{13}

\scriptexercise{ranking-funds}

\ask In which part does the order matter? \emph{(In (a) and (c); in (b) it
	does not.)}

(a) Order matters, no fund appears twice: $10\cdot9\cdot8 = 720$ podiums.

(b) Each top three appears as a podium in $3! = 6$ orders, so
$720/3! = \binom{10}{3} = 120$.

(c) The full league table is an arrangement of all 10 funds:
$10! = 3\,628\,800$.

\say Part (b) is the toolbox remark in action: each set comes in $k!$
orders.

\begin{bg}{mutual funds, fund ratings and league tables}
	\emph{Mutual fund}: a pool of money from many investors, managed by a
	professional who buys a portfolio of stocks and/or bonds. Investors want to
	know which funds are good, so companies such as Morningstar or Lipper
	\emph{rate and rank} funds, typically by past return adjusted for risk
	(e.g.\ the \emph{Sharpe ratio}: average excess return over the risk-free
	rate divided by the volatility). A \emph{league table} is simply the
	ranking of all funds. The ordered podium (who is first) versus the
	unordered ``top three'' is exactly the permutation/combination distinction:
	a ``top quartile'' badge is a combination, a ``number 1 fund'' title is a
	permutation.

	Do not confuse these \emph{fund} ratings with \emph{credit} ratings of
	bonds (Moody's, S\&P, Fitch), which appear in Exercise~\exref{rating-bonds}.

	\terms mutual fund; fund manager; risk-adjusted return; Sharpe ratio;
	league table; fund rating vs.\ credit rating.

	\sq{Is being number one last year a good reason to buy a fund?}
	\ans Mostly no: studies find that past winners rarely stay on top; a large
	part of the ranking is luck. With many funds, \emph{some} fund will be on
	top by chance --- a counting argument in disguise.
\end{bg}

\subsection{Exercise~\exref{committees} --- Committees with a conflict (10 min)}\slide{14}

\scriptexercise{committees}

\begin{cue}
	Give the class two or three minutes to try (a) and (b) in pairs before
	you show the solution; most of them need the time, and (c) is where they
	go wrong.
\end{cue}

\ask Does the order matter in a committee? \emph{(No --- combinations.)}

Without restrictions: choose the men in $\binom63 = 20$ ways, the women in
$\binom83 = 56$ ways, total $20\cdot56 = 1120$. In each part we subtract the
bad committees.

(a) Bad male trios contain both feuding men; the third man comes from the
other 4: $\binom41 = 4$. Good male trios: $20-4 = 16$, and the total is
$16\cdot56 = 896$.

(b) Similarly for the women: both feuding women in, plus 1 of the other 6:
$6$ bad trios. Good: $56-6 = 50$. Total: $20\cdot 50 = 1000$.

(c) Here the conflict is across the two groups, so we subtract from the grand
total. Bad committees: that man is in, so we need 2 more men from the other 5:
$\binom52 = 10$; that woman is in, so 2 more women from the other 7:
$\binom72 = 21$. Bad committees: $10\cdot21 = 210$. Answer:
$1120 - 210 = 910$.

\subsection{Exercise~\exref{portfolio} --- Choosing a portfolio, finance (7 min)}\slide{15}

\scriptexercise{portfolio}

(a) Equally weighted, so only the \emph{set} of stocks matters: unordered, no
repetition,
\[ \binom{30}{5} = \frac{30\cdot29\cdot28\cdot27\cdot26}{5!} = 142\,506. \]

(b) \ask What is the complement of ``at least one''? \emph{(``None''.)}
Portfolios with no tech stock choose all 5 from the 22 non-tech stocks:
$\binom{22}{5} = 26\,334$. Hence $142\,506 - 26\,334 = 116\,172$.

(c) Choose the tech part and the non-tech part separately, then multiply:
2 of the 8 tech stocks and 3 of the 22 others,
\[ \binom82\binom{22}3 = 28\cdot1540 = 43\,120. \]

\say In the Hungarian lottery, ötöslottó, 5 numbers are drawn from 90, and
the order is irrelevant. The number of possible tickets is
$\binom{90}{5} = 43\,949\,268$. So one ticket wins the jackpot in 1 of about
44 million cases.

\begin{cue}
	Optional remark: with 500 stocks and 20 chosen,
	$\binom{500}{20}\approx 2.7\cdot10^{35}$ --- trying all portfolios is
	impossible, which is why portfolio selection uses optimization, not
	enumeration.
\end{cue}

\begin{bg}{portfolio selection, diversification, and the lottery}
	\emph{Index}: a list of stocks whose average price is tracked (e.g.\ the
	Dow Jones Industrial Average has 30 stocks, the S\&P 500 has about 500).
	\emph{Equally weighted portfolio}: the same amount of money in each chosen
	stock, so only the \emph{set} matters --- a combination.

	\emph{Diversification}: holding several stocks reduces risk, because their
	random fluctuations partly cancel. Engineering analogy: averaging several
	noisy, not perfectly correlated sensors reduces the noise. The part of the
	risk that is specific to a single company (\emph{idiosyncratic risk})
	diversifies away; the part common to all stocks (\emph{market} or
	\emph{systematic risk}) does not. Part (b) and (c) are typical
	\emph{sector constraints}: a fund's rules may require at least one, or at
	most two, stocks from a sector, to avoid concentration (all eggs in the
	tech basket).

	\emph{Portfolio selection} (Markowitz, 1952, Nobel prize 1990): choose
	weights to minimize variance for a given expected return. With a limit on
	the number of stocks (``pick 20 of 500'') this becomes a combinatorial
	(mixed-integer) optimization problem, and the $2.7\cdot 10^{35}$
	candidates explain why brute force is hopeless.

	\emph{Lottery as a negative-expected-value bet.} The organizer pays out
	only part of the ticket revenue as prizes (for state lotteries typically
	around half) and keeps the rest. So the expected payoff of a ticket is
	roughly half its price: on average you lose. Buying a ticket is the
	opposite of an investment; finance calls such a bet ``unfair''. People
	still buy because a tiny chance of a life-changing gain is attractive
	(behavioral finance: people overweight small probabilities).

	\terms idiosyncratic vs.\ systematic risk; diversification; mean--variance
	optimization; expected value; fair vs.\ unfair bet.

	\sq{Why would anyone use equal weights instead of optimizing?}
	\ans Optimized weights depend on estimated means and covariances, which
	are very noisy; the simple $1/N$ rule is surprisingly hard to beat in
	practice.

	\sq{Can a lottery ever be a good bet?}
	\ans In theory, when a jackpot has rolled over many times its expected
	value can exceed the price, but the jackpot is usually shared with other
	winners and taxed, so in practice it hardly ever is.
\end{bg}

\subsection{Exercise~\exref{dance-class} --- Dance class: choose, then pair (3 min, optional)}\slide{16}

\scriptexercise{dance-class}

Three independent steps, so we multiply. First \emph{who}: choose the women,
$\binom{10}{5} = 252$; choose the men, $\binom{12}{5} = 792$. Then
\emph{how}: line up the 5 chosen women; the first can get any of 5 men, the
second any of the remaining 4, and so on: $5! = 120$. Altogether
\[ \binom{10}{5}\binom{12}{5}\,5! = 252\cdot792\cdot120 = 23\,950\,080. \]

\say Lesson: combinations answer ``who?'', permutations answer ``how are they
arranged?'', and we multiply.

\subsection{Exercise~\exref{lattice-paths} --- Lattice paths (9 min)}\slide{17}

\scriptexercise{lattice-paths}

\begin{center}
	\begin{tikzpicture}[scale=.6]
		\draw[gray!50] (0,0) grid (4,3);
		\foreach \x in {0,...,4} \foreach \y in {0,...,3} \fill[navy] (\x,\y) circle (2pt);
		\node[below left] at (0,0) {$A$}; \node[above right] at (4,3) {$B$};
		\draw[teal,very thick] (2,1) circle (6pt);
		\draw[gold,ultra thick,->] (0,0)--(1,0)--(1,1)--(2,1)--(3,1)--(3,3)--(4,3);
	\end{tikzpicture}
\end{center}

\begin{cue}
	Let them count a few paths by hand first: one or two minutes of trying is
	what makes the encoding as a word click.
\end{cue}

\say Every path consists of exactly 7 moves: 4 right steps and 3 up steps, in
some order. The path in the picture is R\,U\,R\,R\,U\,U\,R. So a path is a
word with 4 R's and 3 U's, and we already know how to count those.

(a) The path is fixed once we choose \emph{which} 3 of the 7 moves are
``up'': $\binom73 = 35$.

(b) Split the path at $(2,1)$. From $A$ to $(2,1)$: 3 moves, 1 of them U,
$\binom31 = 3$ paths. From $(2,1)$ to $B$: 4 moves, 2 of them U,
$\binom42 = 6$ paths. Every first half combines with every second half, so
by the counting principle $3\cdot6 = 18$ paths.

\subsection{Summary (4 min)}\slide{18}

\begin{itemize}
	\item Ordered, no repetition: $n!/(n-k)!$; ordered with repetition: $n^k$.
	\item Unordered, no repetition: $\binom nk$ (e.g.\ lattice paths);
	      identical objects: divide by $n_1!\cdots n_r!$.
	\item ``Together'' $\to$ glue into a block; ``not together'' $\to$
	      complement.
\end{itemize}

\say For self-study after class, the study guide on the course website has
four extra problems on today's topics, Exercises~\exref{boys-and-girls}
to~\exref{inviting-friends} --- you do not have to hand them in, but they
are very good preparation for next week.

\subsubsection*{Extra practice in the study guide, Session 1 (with detailed
	answers for you)}
\begin{enumerate}[label=\textbf{E\arabic*.}]
	\item[\textbf{\exref{boys-and-girls}}] 3 boys and 3 girls sit in a row. (a) No restriction: 6 distinct
	      people, $6! = 720$. (b) Boys together and girls together: two
	      blocks, 2 orders of the blocks, $3!$ inside each:
	      $2\cdot3!\cdot3! = 72$. (c) Only the boys must sit together: one
	      block of boys plus 3 girls $=4$ units, $4!$ orders, $3!$ inside the
	      block: $4!\cdot3! = 144$. (d) No two of the same sex adjacent: the
	      pattern is BGBGBG or GBGBGB (2 ways), then $3!$ for the boys and
	      $3!$ for the girls: $2\cdot3!\cdot3! = 72$.
	\item[\textbf{\exref{eight-in-a-row}}] 8 people in a row. (a) $A$ and $B$ next to each other: glue them,
	      7 units, $2\cdot7! = 10\,080$. (b) 4 men and 4 women, no two men or
	      two women adjacent: 2 patterns, $2\cdot4!\cdot4! = 1152$. (c) 4
	      couples, each together: 4 blocks, $4!$ orders, 2 orders inside each
	      block: $4!\cdot2^4 = 384$.
	\item[\textbf{\exref{selling-books}}] 2 books out of 6 math, 7 science and 4 economics books. (a) Same
	      subject:
	      \[ \binom62+\binom72+\binom42 = 15+21+6 = 42. \]
	      (b) Different
	      subjects: choose the pair of subjects and one book of each:
	      \[ 6\cdot7+6\cdot4+7\cdot4 = 42+24+28 = 94. \]
	      Check: every pair of
	      books is in exactly one case, $42+94 = 136 = \binom{17}{2}$.
	\item[\textbf{\exref{inviting-friends}}] Invite 5 of 8 friends. (a) Two are feuding and won't come together:
	      all invitations minus those with both:
	      \[ \binom85-\binom63 = 56-20 = 36. \]
	      (b) Two only come together: both invited (3 more from 6),
	      $\binom63 = 20$, or both left out (5 from the other 6), $\binom65=6$;
	      total $26$.
\end{enumerate}

\begin{cue}
	Slide 18 ends with ``Short break --- Session 2 follows right after.''
	Announce the exact restart time (e.g.\ ``we continue in ten minutes'') and
	leave the toolbox table on the board.
\end{cue}

%======================================================================
\section{Session 2: Multinomial coefficients, sample space, events}
%======================================================================

\subsection{After the break: a two-minute recap (3 min)}\slide{19--20}

Welcome back. Let's continue where we stopped before the break. Two minutes
of recap, then something new.

\say Every counting problem starts with two questions: does order matter, and
is repetition allowed?

\board Point at the table from Session 1 and fill in the empty cell:
\begin{center}
	\begin{tabular}{lcc}\toprule
		              & ordered     & unordered                \\\midrule
		no repetition & $n!/(n-k)!$ & $\binom nk$              \\
		repetition    & $n^k$       & $\binom{n+k-1}{k}$ (now) \\\bottomrule
	\end{tabular}
\end{center}

Identical objects in a row: $n!/(n_1!\cdots n_r!)$. ``Together'' means a
block; ``not together'' or ``at least one'' means the complement. And
lattice paths: choose which of the $n$ moves go up, $\binom nk$.

\say Now: we divide objects into several groups, we put identical objects
into boxes, and then we learn the language of events.

\subsection{Toolbox (5 min)}\slide{21}

\emph{Multinomial coefficient.} Suppose we split $n$ distinct objects into
$r$ labeled groups of sizes $n_1,\dots,n_r$ (with $n_1+\dots+n_r=n$). We
choose the first group in $\binom{n}{n_1}$ ways, the second from the
remaining objects in $\binom{n-n_1}{n_2}$ ways, and so on. Multiplying and
canceling gives
\[ \binom{n}{n_1,n_2,\dots,n_r} = \frac{n!}{n_1!\,n_2!\cdots n_r!}. \]
This is the same number as the permutations with repetition: label every
object with the name of its group, and you get a word with $n_1$ copies of
the first label, $n_2$ copies of the second, \ldots

\emph{Binomial and multinomial theorems.}
\[ (a+b)^n = \sum_{k=0}^{n}\binom nk a^kb^{n-k},\qquad
	(x_1+\dots+x_r)^n = \sum_{n_1+\dots+n_r=n}\binom{n}{n_1,\dots,n_r}
	x_1^{n_1}\cdots x_r^{n_r}. \]
\say Why is the coefficient of $a^kb^{n-k}$ equal to $\binom nk$? When you
multiply out $(a+b)(a+b)\cdots(a+b)$, you choose $a$ or $b$ from every
factor. To get $a^kb^{n-k}$ you must choose $a$ from exactly $k$ factors ---
there are $\binom nk$ ways to choose which ones. That is why the numbers
$\binom nk$ are called \emph{binomial coefficients}.

\emph{Stars and bars.} Put $n$ \emph{identical} objects into $r$ distinct
boxes. Draw the objects as $n$ stars and separate the boxes by $r-1$ bars.
Every distribution is one arrangement of $n$ stars and $r-1$ bars, so the
number of distributions is
\[ \binom{n+r-1}{r-1}. \]
If every box must be nonempty, put one object in each box first and
distribute the remaining $n-r$ freely:
\[ \binom{(n-r)+r-1}{r-1} = \binom{n-1}{r-1}. \]

\subsection{Exercise~\exref{bridge} --- Bridge: two correct solutions (5 min)}\slide{22}

\scriptexercise{bridge}

\say I will show you two solutions that look completely different, and both
are correct.

\emph{Permutation view.} Line up the deck: $52!$ orders. Give cards 1--13 to
the first player, 14--26 to the second, and so on. But the order of cards
inside a hand does not matter to the player, so the 13 cards of each hand
behave like identical objects. We divide by $13!$ four times:
\[ \frac{52!}{13!\,13!\,13!\,13!}. \]

\emph{Combination view.} Player 1 gets 13 of the 52 cards, player 2 gets 13
of the remaining 39, player 3 gets 13 of the remaining 26, and player 4 gets
the last 13 (one way):
\[ \binom{52}{13}\binom{39}{13}\binom{26}{13}\binom{13}{13}. \]

\ask Are these equal? Write out the binomial coefficients:
\[ \frac{52!}{13!\,39!}\cdot\frac{39!}{13!\,26!}\cdot\frac{26!}{13!\,13!}
	\cdot 1 = \frac{52!}{(13!)^4}. \]
The $39!$ and the $26!$ cancel telescopically, so both views give the same
number, about $5.36\cdot10^{28}$.

\subsection{Exercise~\exref{dividing-into-groups} --- Dividing into groups (4 min)}\slide{23}

\scriptexercise{dividing-into-groups}

(a) Fill the committees one after the other:
\[ \binom{12}{3}\binom{9}{4}\binom55 = 220\cdot126\cdot1
	= \frac{12!}{3!\,4!\,5!} = 27\,720. \]

(b) Each teacher independently chooses one of 4 schools: $4^8 = 65\,536$.
\say Careful --- this part is not a multinomial problem, because the group
sizes are not fixed. With exactly 2 per school it is: groups of sizes
$2,2,2,2$,
\[ \binom82\binom62\binom42\binom22 = 28\cdot15\cdot6\cdot1
	= \frac{8!}{2!\,2!\,2!\,2!} = 2520. \]

\subsection{Exercise~\exref{rating-bonds} --- Rating bonds, finance (2 min)}\slide{24}

\scriptexercise{rating-bonds}

\ask What are the groups here? \emph{(The three rating classes, of sizes 4,
	3 and 3.)}

Three labeled groups of sizes 4, 3, 3:
\[ \binom{10}{4,3,3} = \frac{10!}{4!\,3!\,3!} = \binom{10}{4}\binom63
	= 210\cdot 20 = 4200. \]
(After choosing the 4 A's, the 3 B's are chosen from the remaining 6, and
the last 3 bonds get C.)

\begin{bg}{credit ratings and rating agencies}
	\emph{Credit rating}: an opinion about how likely a borrower is to pay its
	debts. The three big \emph{rating agencies} are S\&P, Moody's and Fitch.
	Their letter scales run from AAA (S\&P/Fitch) or Aaa (Moody's), the
	safest, through AA, A, BBB, down to BB, B, CCC, \ldots\ and D (default).
	BBB$-$ and above is \emph{investment grade}; below that is \emph{high
		yield}, popularly ``junk''. Many institutional investors (pension funds,
	insurers) may only hold investment-grade bonds, so a downgrade below BBB$-$
	forces selling. Banks also run \emph{internal} rating systems for their
	loan clients; under the Basel rules their capital requirements depend on
	these grades.

	The exercise fixes how many bonds get each grade: that is like a
	\emph{forced distribution} (grading on a curve), used e.g.\ when a bank
	wants its internal grades to match a target distribution. Each rating
	assignment is a labeled division into groups --- a multinomial
	coefficient.

	\terms investment grade vs.\ high yield; downgrade; default probability;
	rating migration (moving between grades over time).

	\sq{Does rating A mean the bond cannot default?}
	\ans No. It means a low estimated default probability (historically,
	A-rated corporate bonds default rarely within a year, but not never). In
	2008 many AAA-rated mortgage securities suffered heavy losses.

	\sq{Who pays the rating agency?}
	\ans Usually the issuer of the bond (``issuer pays''), which is a known
	conflict of interest.
\end{bg}

\subsection{Exercise~\exref{expansions} --- Expansions; identities (6 min)}\slide{25}

\scriptexercise{expansions}

(a) Binomial theorem with $a = 3x^2$, $b = y$: the $k$-th term is
$\binom5k 3^kx^{2k}y^{5-k}$. The coefficients $\binom5k$ are
$1,5,10,10,5,1$; multiply by $3^k = 1,3,9,27,81,243$:
\[ 1\cdot1,\ 5\cdot3,\ 10\cdot9,\ 10\cdot27,\ 5\cdot81,\ 1\cdot243
	\;=\; 1,\ 15,\ 90,\ 270,\ 405,\ 243, \]
so (writing the highest power of $x$ first)
\[ (3x^2+y)^5 = 243x^{10} + 405x^8y + 270x^6y^2 + 90x^4y^3 + 15x^2y^4 + y^5. \]

(b) Out of the 4 factors we need $x_1$ once, $(2x_2)$ twice and $(3x_3)$
once. The number of ways to decide which factor gives what is the
multinomial coefficient
\[ \binom{4}{1,2,1} = \frac{4!}{1!\,2!\,1!} = 12, \]
and the constants contribute $1\cdot2^2\cdot3 = 12$. Coefficient:
$12\cdot12 = 144$.

(c) \say To add up all coefficients, set every variable equal to 1:
$(1+2+3)^4 = 6^4 = 1296$.

\say The same substitution trick gives three useful identities --- you will
find them again in Exercise~\exref{binomial-sums}:
\[ \sum_k\binom nk = (1+1)^n = 2^n,\quad
	\sum_k\binom nk2^k = (2+1)^n = 3^n,\quad
	\sum_k(-1)^k\binom nk = (1-1)^n = 0. \]
The first says: a set with $n$ elements has $2^n$ subsets (each element is
in or out). We will meet it again as the row sum of Pascal's triangle.

\subsection{Exercise~\exref{stars-and-bars} --- Stars and bars (4 min)}\slide{26}

\scriptexercise{stars-and-bars}

(a) Encode a division as 8 stars and 3 bars; the bars separate the 4
schools.
\board \texttt{**|***||***} means $(2,3,0,3)$: 2 boards to school 1, 3 to
school 2, none to school 3, 3 to school 4.

Every division gives exactly one code and every code gives one division. A
code has $8+3 = 11$ symbols and is fixed once we choose the positions of the
3 bars: $\binom{11}{3} = 165$.

(b) At least one each: first give every school one board, then divide the
remaining $4$ freely: 4 stars, 3 bars, $\binom{4+3}{3}=\binom{7}{3} = 35$.

\subsection{Exercise~\exref{identical-dice} --- Two identical dice (2 min)}\slide{27}

\scriptexercise{identical-dice}

Since we cannot tell the dice apart, $(1,2)$ and $(2,1)$ look the same.
Outcomes with two different numbers: $\binom62 = 15$. Doubles: 6. Total 21.

\say Alternatively: we choose 2 numbers out of 6, order does not matter,
repetition is allowed. That is stars and bars again: 2 stars distributed
among the 6 numbers, $\binom{6+2-1}{2} = \binom72 = 21$.

\begin{cue}
	Warn them: the 21 outcomes are \emph{not} equally likely (a double has
	probability $1/36$, a mixed pair $2/36$). This is an important point for
	next week.
\end{cue}

\subsection{Exercise~\exref{allocating-capital} --- Allocating capital, finance (5 min)}\slide{28}

\scriptexercise{allocating-capital}

(a) First pay the minimums: $2+2+3+4 = 11$ units. The remaining $9$ identical
units can be distributed freely among the 4 opportunities --- 9 stars, 3
bars: $\binom{9+3}{3}=\binom{12}{3} = 220$.

(b) Add the strategies in which exactly one opportunity gets nothing. If the
first is left out, we pay the other three minimums, $2+3+4 = 9$, and
distribute the remaining 11 units among 3 opportunities (11 stars, 2 bars):
$\binom{13}{2} = 78$. The other cases are in the table:

\begin{center}
	\begin{tabular}{lccc}\toprule
		left out    & minimums paid & left to divide & count                 \\\midrule
		1st (min 2) & $2+3+4=9$     & 11 among 3     & $\binom{13}{2} = 78$  \\
		2nd (min 2) & $2+3+4=9$     & 11 among 3     & $78$                  \\
		3rd (min 3) & $2+2+4=8$     & 12 among 3     & $\binom{14}{2} = 91$  \\
		4th (min 4) & $2+2+3=7$     & 13 among 3     & $\binom{15}{2} = 105$ \\\bottomrule
	\end{tabular}
\end{center}

Total:
\[ 220+78+78+91+105 = 572. \]

\begin{cue}
	If asked why we cannot leave out two opportunities: (b) says ``at least
	3 of the 4''. Leaving out two would be a different question.
\end{cue}

\begin{bg}{capital allocation with minimum investments}
	\emph{Capital allocation}: deciding how to split a budget among
	investment opportunities. Real opportunities often come in \emph{lots} and
	have \emph{minimum tickets}: a private-equity fund may accept commitments
	only above a threshold, some bonds trade in minimum denominations (e.g.\
	100\,000 euros), and many funds have a minimum subscription. That is why
	the amounts are integers and each opportunity is either skipped or gets at
	least its minimum.

	Counting the strategies tells you the size of the search space: 572
	strategies is tiny, a computer can evaluate every one of them (brute
	force). With 50 opportunities and a budget of 1000 units the count would
	be astronomically large, and one uses \emph{integer programming} instead
	--- the same algorithms engineers use for scheduling and network design.

	\terms capital allocation; minimum ticket (minimum subscription); lot
	size; integer programming; search space.

	\sq{Why not invest fractional amounts?}
	\ans Then the set of strategies is continuous (a simplex) and one uses
	ordinary optimization; the integer version is what you face with lots and
	minimum tickets.

	\sq{Which of the 572 strategies is best?}
	\ans Counting cannot say; it depends on expected returns and risks. That
	is the portfolio-selection problem from Exercise~\exref{portfolio}.
\end{bg}

\subsection{Exercise~\exref{pascals-rule} --- A combinatorial proof (3 min)}\slide{29}

\scriptexercise{pascals-rule}

\say After the binomial tree, this rule will build Pascal's triangle. Here
are two proofs.

\emph{Algebra} \emph{(optional if you are short of time --- it is on the
	slide).} Bring both fractions to the common denominator $(k+1)!\,(n-k)!$:
\[ \frac{n!}{k!(n-k)!}+\frac{n!}{(k+1)!(n-k-1)!}
	= \frac{n!\,\big[(k+1)+(n-k)\big]}{(k+1)!\,(n-k)!}
	= \frac{(n+1)!}{(k+1)!\,(n-k)!} = \binom{n+1}{k+1}, \]
because $n!\,(n+1) = (n+1)!$ and $n-k = (n+1)-(k+1)$.

\emph{Story.} \say Here is a proof with no algebra at all. A bowl contains
$n$ apples and 1 pear. In how many ways can we take $k+1$ fruits? Obviously
$\binom{n+1}{k+1}$. Now split the choices according to the pear. If the pear
is not taken, all $k+1$ fruits are apples: $\binom n{k+1}$. If the pear is
taken, we need $k$ more apples from $n$: $\binom nk$. We counted the same set
in two ways, so the two answers are equal.

\say For self-study, Exercise~\exref{vandermonde} in the study guide asks you to
prove
\[ \sum_{j=0}^k\binom nj\binom m{k-j} = \binom{n+m}k \]
in the same way.
\emph{(Hint, if asked: $n$ apples and $m$ pears, choose $k$ fruits, split by
	the number $j$ of apples.)}

\subsection{Exercise~\exref{binomial-tree} --- The binomial tree, finance (5 min)}\slide{30}

\scriptexercise{binomial-tree}

\begin{cue}
	Point at the small tree on the slide: the number in each node is the
	number of paths reaching it, and each node is the sum of the two nodes
	feeding it --- Pascal's rule, which we have just proved. The same
	lattice-path idea as in Exercise~\exref{lattice-paths}, rotated by 45 degrees. The next
	slide arranges these numbers in Pascal's triangle.
\end{cue}

\say Before the break we counted lattice paths; here is the same counting
in finance.

(a) Two choices on each of 10 days: $2^{10} = 1024$ paths.

(b) A path is a word of 6 U's and 4 D's; choose the 6 up-days:
$\binom{10}{6} = 210$.

(c) The order of the factors does not change the product ($ud = du$), so the
final price depends only on $k\in\{0,1,\dots,10\}$: 11 final prices. The
tree \emph{recombines}: ``up then down'' gives the same price as ``down then
up''.

\say Check: grouping the paths by $k$ gives all paths,
\[ \sum_{k=0}^{10}\binom{10}{k} = 2^{10}. \]
This is exactly the combinatorics
behind the Cox--Ross--Rubinstein option-pricing model you will meet in
derivatives courses.

\begin{bg}{the binomial tree and the binomial option-pricing model}
	\emph{Option.} A \emph{call option} with strike $K$ and maturity $n$ gives
	the right to buy the stock at price $K$ at time $n$. If the stock then
	costs $S_n$, the call pays $\max(S_n-K,0)$ (you buy at $K$ and can sell at
	$S_n$; if $S_n<K$ you simply do not use the right). A \emph{put} pays
	$\max(K-S_n,0)$. The question is: what is a fair price \emph{today}?

	\emph{The model} (Cox, Ross and Rubinstein, 1979). Time is split into $n$
	steps; in each step the price is multiplied by $u>1$ or $d<1$. Engineering
	view: a discretisation of a continuous random process (geometric Brownian
	motion), much like a finite-difference grid; as $n\to\infty$ the price
	converges to the Black--Scholes formula. Typically
	$u = e^{\sigma\sqrt{\Delta t}}$, $d=1/u$, where $\sigma$ is the
	volatility.

	\emph{Pricing by no-arbitrage.} One can build a portfolio of stock and
	cash that pays exactly what the option pays in both successor nodes
	(\emph{replication}); the option must cost what that portfolio costs, or
	someone could make riskless profit (\emph{arbitrage}). The result can be
	written as an expectation with an artificial ``risk-neutral'' probability
	\[ q = \frac{(1+r)-d}{u-d} \]
	of an up-step ($r$ = risk-free interest per step):
	\[ C_0 = \frac{1}{(1+r)^n}\sum_{k=0}^n \binom nk q^k(1-q)^{n-k}
		\max\big(S_0u^kd^{n-k}-K,\,0\big). \]
	\emph{Why paths are counted}: the payoff depends only on the final node
	$k$, and every one of the $\binom nk$ paths into node $k$ has the same
	probability $q^k(1-q)^{n-k}$. So the binomial coefficient is the number of
	ways to get there. The recombining tree has only $n+1$ final nodes instead
	of $2^n$ leaves --- that is what makes the model computable.

	\emph{Tiny example.} $S_0=100$, $u=1.1$, $d=0.9$, $r=0$, $n=2$, $K=100$:
	$q=0.5$; final prices $121, 99, 81$ with $1,2,1$ paths; the call pays
	$21,0,0$; price $= \frac14\cdot21 = 5.25$.

	\terms call/put; strike; maturity; volatility; replication; arbitrage;
	risk-neutral probability; recombining tree.

	\sq{Why not use the real probability of an up-move?}
	\ans Because the price comes from replication, not from a bet on the
	direction: two investors who disagree about the real probability must
	still agree on the cost of the replicating portfolio. The real probability
	does not appear in the formula.

	\sq{Are daily up/down moves of fixed size realistic?}
	\ans Not for a single day, but with many small steps the tree approximates
	a continuous price process very well; it is a numerical method, like a
	grid in engineering.
\end{bg}

\subsection{Pascal's triangle (4 min)}\slide{31}

\say The numbers in the binomial tree are worth a closer look. The number of
paths with $n$ moves, $k$ of them up, is $\binom nk$. Write these numbers in
rows: row $n$ contains $\binom n0,\binom n1,\dots,\binom nn$. This is
\emph{Pascal's triangle}.

\begin{center}
	\begin{tikzpicture}[every node/.style={font=\small}]
		\foreach \n in {0,...,6} {
				\node[navy!70] at (-4.2,-.62*\n) {\scriptsize $n=\n$};
				\foreach \k in {0,...,\n} {
						\pgfmathtruncatemacro{\c}{factorial(\n)/(factorial(\k)*factorial(\n-\k))}
						\node[inner sep=1pt,minimum size=13pt] (p\n-\k)
						at ({(\k-\n/2)*.92},{-.62*\n}) {\c};
					}
			}
		\draw[gold,thick] (p4-1) circle (7pt);
		\draw[gold,thick] (p4-2) circle (7pt);
		\draw[teal,thick,fill=teal!15,fill opacity=.5] (p5-2) circle (7pt);
		\draw[gold,thick,->] (p4-1)--(p5-2);
		\draw[gold,thick,->] (p4-2)--(p5-2);
		\node[teal] at (p5-2) {\small 10};
	\end{tikzpicture}
\end{center}

\board Draw the first rows of the triangle (at least up to $n=5$) and
\[ \binom nk+\binom n{k+1}=\binom{n+1}{k+1},\qquad\text{e.g. }
	\binom41+\binom42 = 4+6 = 10 = \binom52. \]

\say The rule we just proved in Exercise~\exref{pascals-rule} builds the triangle: each
entry is the sum of the two entries above it. You saw the same thing in the
binomial tree: the number of paths into a node is the sum of the paths into
the two nodes feeding it. And in the lattice of Exercise~\exref{lattice-paths}: to reach a
point you arrive either from the point on its left or from the point below
it.

\ask Who can tell me the next row, $n=7$? \emph{(1, 7, 21, 35, 35, 21, 7, 1
	--- each entry the sum of the two above it.)}

\say Ways to read row $n$:
\begin{itemize}
	\item paths: $\binom nk$ paths with $k$ up-moves among $n$ moves --- the
	      lattice paths of Exercise~\exref{lattice-paths} and the binomial
	      tree of Exercise~\exref{binomial-tree}, just before;
	\item the coefficients of $(a+b)^n$, by the binomial theorem. For example,
	      row 5 gives
	      \[ (a+b)^5 = a^5+5a^4b+10a^3b^2+10a^2b^3+5ab^4+b^5; \]
	\item the sum of row $n$ is $2^n$ (for row 5: $1+5+10+10+5+1 = 32$), and
	      every row is symmetric, $\binom nk = \binom n{n-k}$: choosing the $k$
	      elements you take is the same as choosing the $n-k$ you leave.
\end{itemize}

\begin{bg}{Pascal's triangle in finance}
	The same numbers are the path counts of the binomial tree of a stock
	price (Exercise~\exref{binomial-tree}). In the binomial option model the row-$n$
	entries are the \emph{path multiplicities} of the final nodes; with risk-neutral probability $q$, node
	$k$ gets weight $\binom nk q^k(1-q)^{n-k}$ --- the \emph{binomial
		distribution}, which students will meet in a few weeks. With $q=1/2$ the
	weights are simply row $n$ divided by $2^n$ (e.g.\ for $n=2$:
	$\frac14,\frac24,\frac14$, as in the tiny example in the box of
	Exercise~\exref{binomial-tree}). Practitioners
	build trees with hundreds of steps; the triangle is never written out, but
	the same ``add the two parents'' recursion is how the tree is evaluated
	numerically (\emph{backward induction}).

	\terms path multiplicity; binomial distribution; risk-neutral weight;
	backward induction.

	\sq{Why do the middle nodes matter most in the binomial model?}
	\ans Because most paths end there: in row 10 the middle entry is
	$\binom{10}{5} = 252$ of the $1024$ paths, the two extreme nodes only 1
	each. So (with $q$ near $1/2$) the final price is most likely near the
	middle, and for many steps the weights look like a bell curve.
\end{bg}

%----------------------------------------------------------------------
\subsection{Sample space and events (5 min)}\slide{32}

\say Now we switch to the language of probability.

The \emph{sample space} $\Om$ is the set of all possible outcomes of an
experiment. An \emph{event} is a subset $E\subseteq\Om$. We say that $E$
\emph{occurs} if the outcome of the experiment is an element of $E$.

Because events are sets, we can combine them:
\begin{itemize}
	\item $E\cup F$: $E$ \emph{or} $F$ occurs (at least one of them);
	\item $E\cap F$, often written $EF$: \emph{both} occur;
	\item $E^c$: $E$ does \emph{not} occur;
	\item $E\subseteq F$: whenever $E$ occurs, $F$ occurs too ($E$ implies $F$);
	\item $E$ and $F$ are \emph{mutually exclusive} if $EF=\emptyset$: they
	      cannot occur together.
\end{itemize}

\begin{center}
	\setlength{\tabcolsep}{5mm}
	\begin{tabular}{ccc}
		\venn{\fill[teal!35] (0,0) circle (.85);}
		          & \venn{\fill[teal!35] (1,0) circle (.85);}
		          & \venn{\fill[teal!35] (0,0) circle (.85); \fill[white] (1,0) circle (.85);}                       \\
		$E$       & $F$                                                                        & $E\setminus F=EF^c$ \\[2mm]
		\venn{\fill[teal!35] (0,0) circle (.85); \fill[teal!35] (1,0) circle (.85);}
		          & \venn{\begin{scope}\clip (0,0) circle (.85);
				            \fill[teal!35] (1,0) circle (.85);\end{scope}}
		          & \venn{\shadeall{teal!35}\fill[white] (0,0) circle (.85);}                                        \\
		$E\cup F$ & $EF$                                                                       & $E^c$
	\end{tabular}
\end{center}

\say A Venn diagram is a picture of this: the rectangle is $\Om$, the circles
are the events, and the shaded region is the event we are talking about.
Let's read the six pictures.

\begin{itemize}
	\item Top row: $E$ alone is the left circle, $F$ alone the right circle.
	      The third picture is $E\setminus F$, ``$E$ but not $F$'': the left
	      circle without the lens. In event notation this is $EF^c$ ---
	      $E$ occurs \emph{and} $F$ does not.
	\item Bottom row: $E\cup F$ shades both circles (at least one occurs);
	      $EF$ shades only the lens in the middle (both occur); $E^c$ shades
	      everything outside the left circle ($E$ does not occur).
\end{itemize}

\ask Which picture shows ``$F$ occurs but $E$ does not''? \emph{(None of the
	six: it is $FE^c=F\setminus E$, the right circle without the lens.)}

\subsection{De Morgan's laws for two events (5 min)}\slide{33}

\board
\[ (E\cup F)^c = E^c\cap F^c, \qquad (E\cap F)^c = E^c\cup F^c. \]

\say In words: ``not ($E$ or $F$)'' is the same as ``neither $E$ nor $F$''.
And ``not (both $E$ and $F$)'' is the same as ``at least one of them fails''.

\begin{center}
	\renewcommand{\arraystretch}{1.3}
	\begin{tabular}{ccccc}
		\venn{\shadeall{teal!35}\fill[white] (0,0) circle (.85);}
		                  & $\cap$ &
		\venn{\shadeall{gold!40}\fill[white] (1,0) circle (.85);}
		                  & $=$    &
		\venn{\shadeall{navy!35}\fill[white] (0,0) circle (.85);
			\fill[white] (1,0) circle (.85);}       \\[-1mm]
		\scriptsize $E^c$ &        & \scriptsize $F^c$ &  &
		\scriptsize $E^c\cap F^c$: outside both $=(E\cup F)^c$        \\
		\venn{\shadeall{teal!35}\fill[white] (0,0) circle (.85);}
		                  & $\cup$ &
		\venn{\shadeall{gold!40}\fill[white] (1,0) circle (.85);}
		                  & $=$    &
		\venn{\shadeall{navy!35}\begin{scope}\clip (0,0) circle (.85);
				\fill[white] (1,0) circle (.85);\end{scope}} \\[-1mm]
		\scriptsize $E^c$ &        & \scriptsize $F^c$ &  &
		\scriptsize $E^c\cup F^c$: all but the lens $=(E\cap F)^c$
	\end{tabular}
\end{center}

\say Look at the pictures. Shade $E^c$ --- everything outside $E$ --- in one
color, and $F^c$ in another. The region shaded \emph{twice} is the
intersection $E^c\cap F^c$: that is the part outside both circles, which is
exactly the complement of $E\cup F$. The region shaded \emph{at least once}
is the union $E^c\cup F^c$: that is everything except the lens in the
middle, and the lens is $E\cap F$. So the union is $(E\cap F)^c$.

\ask Give me an everyday example of the second law. \emph{(``It is not true
	that I passed both exams'' = ``I failed at least one of them''.)}

\subsection{De Morgan's laws: general form (2 min)}\slide{34}

\say The same laws hold for any number of events: the complement of a union
is the intersection of the complements, and the complement of an
intersection is the union of the complements.

\board General form, $n$ events:
\[ \Big(\bigcup_{i=1}^n E_i\Big)^c = \bigcap_{i=1}^n E_i^c,\qquad
	\Big(\bigcap_{i=1}^n E_i\Big)^c = \bigcup_{i=1}^n E_i^c. \]

\say Next week we will use exactly this:
$P(A^cB^c) = 1-P(A\cup B)$.

\subsection{Exercise~\exref{sample-space} --- Writing down the sample space (3 min)}\slide{35}

\scriptexercise{sample-space}

\say Record the outcome as an ordered pair: (first draw, second draw).

With replacement: 3 choices, then again 3 choices:
\[ \Om = \{rr,rg,rb,gr,gg,gb,br,bg,bb\},\qquad |\Om| = 3\cdot3 = 9. \]

Without replacement: the second draw cannot repeat the first, so the pairs
$rr, gg, bb$ disappear:
\[ \Om = \{rg,rb,gr,gb,br,bg\},\qquad |\Om| = 3\cdot2 = 6. \]

\subsection{Exercise~\exref{stock-path-events} --- Events of a stock path, finance (4 min)}\slide{36}

\scriptexercise{stock-path-events}

\say Track the position --- the number of ups minus the number of downs ---
after each day. (Up and down steps have the same size here.)

\begin{itemize}
	\item $A$: at least 2 U's, so the net position is $+1$ or $+3$:
	      \[ A = \{UUU,UUD,UDU,DUU\}. \]
	\item $B$: the first letter is U:
	      \[ B = \{UUU,UUD,UDU,UDD\}. \]
	\item $C$: the positions must stay $\ge 0$. $UDD$ goes $+1,0,-1$, so it
	      drops below; any path starting with $D$ goes to $-1$ immediately.
	      Hence
	      \[ C = \{UUU,UUD,UDU\}. \]
\end{itemize}

\ask What is the relation between $C$ and $A$? \emph{(After 3 steps the
	position is odd, so ``never below'' forces it to end at $\ge1$: $C\subseteq
		A$.)}

Further:
\[ \begin{aligned}
		AB   & = \{UUU,UUD,UDU\} = C;                          \\
		A^cB & = \{UDD\}\quad\text{(up first, but ends below)};   \\
		AB^c & = \{DUU\}\quad\text{(down first, but ends above)}.
	\end{aligned} \]

\begin{bg}{path-dependent events and barrier options}
	Event $A$ depends only on the final price; event $C$ depends on the whole
	path. Derivatives whose payoff depends on the path are called
	\emph{path-dependent}. The classic example is a \emph{barrier option}: a
	\emph{down-and-out call} is an ordinary call option that is canceled
	(``knocked out'') if the price ever falls below a barrier. If the barrier
	is the starting price, the option is alive at maturity exactly on event $C$.
	Because it can be knocked out, it is cheaper than the plain call --- you
	give up the scenarios in $A\setminus C=\{DUU\}$ (ends above, but touched the
	barrier). Engineering analogy: a circuit breaker --- it does not matter
	where the current ends up, only whether it ever exceeded the limit.

	\terms path-dependent; barrier; knock-out / knock-in; Asian option (pays
	on the average price).

	\sq{Why would anyone buy a knock-out option?}
	\ans It is cheaper, and an investor who believes the price will not fall
	below the barrier does not want to pay for the scenarios where it does.
\end{bg}

\subsection{Exercise~\exref{two-dice} --- Two dice: event algebra (5 min)}\slide{37}

\scriptexercise{two-dice}

\[ \Om = \{(i,j): i,j = 1,\dots,6\},\qquad |\Om| = 36. \]
The sum is odd exactly when
one die is even and the other odd: 3 even values, 3 odd values, and 2 ways
to decide which die is the even one, so $|E| = 3\cdot3\cdot2 = 18$. The
complement of $F$ is ``no 1 on either die'', which has $5\cdot5 = 25$
outcomes, so $|F| = 36-25 = 11$.

\begin{itemize}
	\item $EF$: one die is 1, which is odd, so the other must be even:
	      \[ \{(1,2),(1,4),(1,6),(2,1),(4,1),(6,1)\}, \]
	      6 outcomes.
	\item $E\cup F$:
	      \[ |E|+|F|-|EF| = 18+11-6 = 23 \]
	      outcomes: the 18 odd sums
	      plus the elements of $F$ with an even sum,
	      \[ (1,1),(1,3),(1,5),(3,1),(5,1). \]
	\item $G = \{(1,4),(2,3),(3,2),(4,1)\}$, so $FG = \{(1,4),(4,1)\}$.
	\item $EF^c$: odd sum and no 1: $|E|-|EF| = 18-6 = 12$ outcomes.
	\item $5$ is odd, so $G\subseteq E$, and $EFG = FG = \{(1,4),(4,1)\}$.
\end{itemize}

\say The count in the second bullet is the inclusion--exclusion formula,
$|E\cup F| = |E|+|F|-|EF|$: in $|E|+|F|$ the overlap was counted twice.
Next week it becomes $P(E\cup F) = P(E)+P(F)-P(EF)$.

\subsection{Exercise~\exref{credit-basket} --- A credit basket, finance (5 min)}\slide{38}

\scriptexercise{credit-basket}

\begin{enumerate}[label=(\alph*)]
	\item $D_1\cup D_2\cup D_3$. This is where a \emph{first-to-default swap}
	      pays. It contains all outcomes except $000$: $8-1 = 7$ outcomes.
	\item $(D_1\cup D_2\cup D_3)^c = D_1^cD_2^cD_3^c$ by De Morgan: only $000$,
	      1 outcome.
	\item $D_1D_2^cD_3^c\ \cup\ D_1^cD_2D_3^c\ \cup\ D_1^cD_2^cD_3$: ``bond 1
	      alone'', ``bond 2 alone'', ``bond 3 alone''. The three pieces are
	      mutually exclusive: $100, 010, 001$, i.e.\ $\binom31 = 3$ outcomes.
	\item $D_1D_2\cup D_1D_3\cup D_2D_3$: some pair defaults. This is where a
	      \emph{second-to-default swap} pays. Outcomes $110, 101, 011, 111$:
	      exactly two defaults, $\binom32=3$, plus all three, $\binom33=1$;
	      4 outcomes.
\end{enumerate}
Check:
\[ \text{(b)} + \text{(c)} + \text{(d)} = 1+3+4 = 8 = |\Om|,\qquad\text{and}\qquad
	\text{(a)} = \text{(c)} + \text{(d)} = 3+4 = 7. \]

\begin{bg}{credit default swaps, first- and second-to-default}
	\emph{Credit default swap (CDS)}: insurance on a bond. The protection
	buyer pays a regular premium (the \emph{CDS spread}, quoted in basis
	points per year); if the reference company defaults, the protection seller
	pays the loss, i.e.\ face value minus what can be recovered (the
	\emph{recovery rate}, often around 40\%).

	\emph{Basket default swap}: protection on a basket of several names
	(typically 3--10). An \emph{$n$-th-to-default swap} pays when the $n$-th
	default in the basket happens, and then terminates.
	\begin{itemize}[nosep]
		\item \emph{First-to-default (FtD)}: pays at the first default, i.e.\ on
		      the event ``at least one default'' $=D_1\cup D_2\cup D_3$.
		\item \emph{Second-to-default (StD)}: pays at the second default, i.e.\
		      on ``at least two defaults'' $=D_1D_2\cup D_1D_3\cup D_2D_3$.
	\end{itemize}
	The FtD is the riskier position for the seller (its event is bigger), so
	its spread is higher. Illustration (next week's material): if each bond
	defaults with probability 2\% independently,
	\[ P(\text{at least one}) = 1-0.98^3 \approx 5.9\%, \]
	while
	\[ P(\text{at least two}) = 3(0.02)^2(0.98)+0.02^3\approx 0.12\%. \]

	\emph{Correlation} matters: if defaults tend to happen together, ``at least
	one'' becomes less likely than under independence (defaults cluster in the
	same scenarios) and ``at least two'' becomes more likely; FtD gets cheaper
	and StD more expensive.

	\emph{Structured products} push the same idea further: a \emph{CDO
		tranche} or a \emph{credit-linked note} pays depending on \emph{which}
	combination of names in a pool survives --- in engineering terms a
	reliability block diagram of series and parallel ``routes''. The lesson
	that transfers best from engineering is \emph{common-cause failure}:
	redundancy helps only if failures are not strongly correlated. In 2008
	mortgage defaults all depended on house prices, so ``redundant''
	structures failed together, like parallel components sharing one power
	supply.

	\terms CDS spread; protection buyer/seller; recovery rate; basket; default
	correlation; structured product; CDO tranche; common-cause failure.

	\sq{Why buy an FtD instead of three separate CDSs?}
	\ans It is cheaper; it covers only the first loss, which is enough if you
	think one default is plausible but several are not.

	\sq{What happens to the swap after the first default?}
	\ans The FtD pays and ends; the StD keeps running and now waits for one
	more default.
\end{bg}

\subsection{Exercise~\exref{newspapers} --- Venn diagrams and inclusion--exclusion (6 min)}\slide{39--40}

\scriptexercise{newspapers}

\say Fill the Venn diagram from the inside out: first the center, then the
regions of exactly two papers, then the regions of exactly one paper.

\begin{center}
	\begin{tikzpicture}[scale=.75]
		\draw[navy,thick] (0,0) circle (1.5);
		\draw[teal,thick] (1.8,0) circle (1.5);
		\draw[gold,thick] (.9,-1.5) circle (1.5);
		\node at (-.7,.3) {1}; \node at (2.5,.3) {19}; \node at (.9,-2.3) {0};
		\node at (.9,.6) {7}; \node at (.1,-.95) {1}; \node at (1.7,-.95) {3};
		\node at (.9,-.45) {1};
		\node[navy] at (-1.4,1.4) {I}; \node[teal] at (3.2,1.4) {II};
		\node[gold] at (.9,-3.3) {III};
	\end{tikzpicture}
\end{center}

\begin{itemize}
	\item Center (all three): $1$.
	\item Only I\&II: $8-1 = 7$; only I\&III: $2-1 = 1$; only II\&III:
	      $4-1 = 3$. (The 8\% reading I and II include the 1\% reading all
	      three.)
	\item Only I: $10-7-1-1 = 1$; only II: $30-7-3-1 = 19$; only III:
	      $5-1-3-1 = 0$.
\end{itemize}

\begin{cue}
	Slide 40 (``(cont.)'') shows the answers; let the class read them off the
	diagram first.
\end{cue}

(a) Exactly one:
\[ \text{only I} + \text{only II} + \text{only III} = 1+19+0 = 20\%, \]
i.e.\ 20\,000 people.

(b) At least two: the four inner regions, $7+1+3+1 = 12\%$: 12\,000.

(c)
\[ (I\cup III)\cap II = (I\cap II)\cup(III\cap II); \]
the two overlap in
``all three'', so $8+4-1 = 11\%$: 11\,000.

(d) Inclusion--exclusion:
\[ |I\cup II\cup III| = 10+30+5-8-2-4+1 = 32\%, \]
so
$100-32 = 68\%$ read no paper: 68\,000.

(e) II together with exactly one of I, III: only I\&II $+$ only II\&III
$= 7+3 = 10\%$: 10\,000.

\board Inclusion--exclusion for three sets:
\[ |E\cup F\cup G| = |E|+|F|+|G|-|EF|-|EG|-|FG|+|EFG|. \]
\say Adding the three sets counts the pairwise overlaps twice, so we subtract
them; but then the center has been added three times and subtracted three
times, so we add it back once.

\subsection{Exercise~\exref{bank-products} --- Bank products, finance (4 min)}\slide{41}

\scriptexercise{bank-products}

(a) Inclusion--exclusion for three sets:
\[ |C\cup K\cup M| = 820+350+240-300-200-90+70 = 890, \]
so $1000-890 = 110$ clients have none of the products.

(b) ``Only $C$'' $= |C|-|CK|-|CM|+|CKM|$: we remove the two overlaps, but
the clients with all three were removed twice, so we add them back once.
Only $C$: $820-300-200+70 = 390$. Only $K$: $350-300-90+70 = 30$. Only $M$:
$240-200-90+70 = 20$. Total: $390+30+20 = 440$.

(c) $C$ splits into $CK$ and $CK^c$, so
\[ |CK^c| = |C| - |CK| = 820-300 = 520. \]

\say Divide every number by 1000 and you have probabilities. That is where we
start next week: the axioms of probability and
\[ P(E\cup F) = P(E)+P(F)-P(EF). \]

\begin{bg}{retail banking and cross-selling}
	\emph{Retail clients} are individuals (as opposed to companies). The three
	products: a \emph{current account} (everyday account for salary and
	payments), a \emph{credit card}, a \emph{mortgage} (a long-term loan for
	buying a home, secured by the home). \emph{Cross-selling} means selling an
	additional product to an existing client. Banks love it because acquiring
	a new client is expensive, while an existing client is already known (the
	bank sees their income on the current account, which lowers credit risk).
	A common bank KPI is ``products per client''.

	Part (c) is a \emph{target list}: clients with $C$ but not $K$ are the
	natural candidates for a card offer --- in set language $C\setminus K =
		CK^c$. The 110 clients with \emph{no} product are probably dormant
	accounts or clients of other services, a different campaign.

	\terms retail client; current account; mortgage; cross-selling; target
	list; products per client (KPI).

	\sq{Is aggressive cross-selling a problem?}
	\ans It can be: in 2016 Wells Fargo was fined after employees, pushed by
	sales targets, opened millions of accounts without customers' consent.
	Regulators now watch sales incentives closely.
\end{bg}

\subsection{Summary (3 min)}\slide{42}

\begin{itemize}
	\item Groups of given sizes: multinomial coefficient; identical objects
	      into boxes: stars and bars.
	\item Identities: count the same set in two ways (Pascal's rule).
	\item Binomial tree: $\binom nk$ of the $2^n$ paths have $k$ up-moves;
	      Pascal's rule arranges these numbers in Pascal's triangle.
	\item Events are sets: $\cup$ = or, $\cap$ = and, $^c$ = not. De Morgan:
	      $(E\cup F)^c = E^cF^c$ and $(EF)^c = E^c\cup F^c$. Fill Venn diagrams
	      from the inside out.
\end{itemize}

\say Next week, in week 9, we turn these counts into probabilities: the
axioms of probability. For self-study, the study guide has extra problems
on today's topics, Exercises~\exref{binomial-sums}
to~\exref{system-reliability}.

\subsubsection*{Extra practice in the study guide, Session 2 (with detailed
	answers for you)}
\begin{enumerate}[label=\textbf{E\arabic*.},before=\setlength{\emergencystretch}{1.5em}]
	\item[\textbf{\exref{binomial-sums}}] $\sum_k\binom nk = 2^n$ and $\sum_k(-1)^k\binom nk = 0$:
	      substitute $a=b=1$, resp.\ $a=-1$, $b=1$ in the binomial theorem:
	      $(1+1)^n = 2^n$ and $(-1+1)^n = 0$. The second identity says that a
	      nonempty set ($n\ge1$) has as many subsets of even size as of odd
	      size.
	\item[\textbf{\exref{vandermonde}}] $\sum_{j=0}^k\binom nj\binom m{k-j} = \binom{n+m}k$: a bowl
	      holds $n$ apples and $m$ pears; choose $k$ fruits. Directly:
	      $\binom{n+m}k$. Split by the number $j$ of apples: $\binom nj$ ways
	      for the apples and $\binom m{k-j}$ for the pears, and add over
	      $j = 0,\dots,k$. Same set counted twice $\Rightarrow$ identity.
	\item[\textbf{\exref{coin-game}}] $A$, $B$, $C$ flip a coin in turn ($A$ first), the first head wins.
	      With 1 = head and 0 = tail,
	      \[ \Om = \{1, 01, 001, 0001, \dots\}\cup\{000\cdots\} \]
	      (the last outcome: tails forever).
	      \begin{itemize}[nosep]
		      \item $A$ wins: $\{1, 0001, 0000001,\dots\}$ (first head on flip
		            $1, 4, 7,\dots$);
		      \item $B$ wins: $\{01, 00001, \dots\}$ (first head on flip
		            $2, 5, 8, \dots$);
		      \item $(A\cup B)^c$: $C$ wins (first head on flip $3,6,9,\dots$)
		            or nobody ever wins ($000\cdots$).
	      \end{itemize}
	\item[\textbf{\exref{team-classification}}] 15 team members, each blue/white collar and R/D/I: each person has
	      $2\cdot3 = 6$ types, so $|\Om| = 6^{15} \approx 4.7\cdot10^{11}$.
	      At least one blue-collar: complement of ``all white-collar'' ($3$
	      types each): $6^{15}-3^{15}$. No Independent: $2\cdot2 = 4$ types
	      each, $4^{15}$.
\end{enumerate}

\begin{cue}
	Week 9 bridge: start with Exercise~\exref{newspapers} and
	Exercise~\exref{bank-products} in probability form, then the week 9
	exercises on the axioms of probability.
\end{cue}

\end{document}
